% (c) Nikita Lisitsa, lisyarus@gmail.com, 2026

\documentclass[10pt]{beamer}

\usepackage[T2A]{fontenc}
\usepackage[russian]{babel}
\usepackage{minted}

\usepackage[most]{tcolorbox}
\tcbuselibrary{minted}

\usepackage{graphicx}
\graphicspath{ {./images/} }

\usepackage{adjustbox}

\usepackage{color}
\usepackage{soul}
\usepackage{multicol}
\usepackage{hyperref}
\usepackage{tikz}
\usetikzlibrary{arrows.meta,positioning,calc}

\usetheme{metropolis}

\definecolor{red}{rgb}{1,0,0}
\definecolor{codebg}{RGB}{29,35,49}
\setminted{fontsize=\footnotesize}

\makeatletter
\newcommand{\slideimage}[1]{
  \begin{figure}
    \begin{adjustbox}{width=\textwidth, totalheight=\textheight-2\baselineskip-2\baselineskip,keepaspectratio}
      \includegraphics{#1}
    \end{adjustbox}
  \end{figure}
}
\makeatother

\title{Язык программирования C}
\subtitle{Лекция 1: языки высокого уровня, программы из нескольких файлов}
\author{Лисица Никита}
\date{2026}

\setbeamertemplate{footline}[frame number]

\usemintedstyle{tango}

\begin{document}

\frame{\titlepage}

\begin{frame}[fragile]
\frametitle{Напоминание из курса информатики :)}
\begin{itemize}
\item Всё, с чем умеет работать компьютер -- это числа:
\pause
\begin{itemize}
\item Текст -- последовательность чисел, интерпретируемая посредством кодировки (напр. в ASCII символ 'A' это число 65)
\pause
\item Картинка -- последовательность чисел, задающих цвета (RGB -- красный, зелёный, синий)
\pause
\item Звук -- последовательность чисел, задающих звуковую волну (часто 16-битные, от -32768 до 32767)
\pause
\item Программа -- последовательность чисел, задающих команды для исполнения процессору
\end{itemize}
\pause
\item Числа хранятся в двоичной системе счисления (есть ток / нет тока)
\end{itemize}
\end{frame}

\begin{frame}[fragile]
\frametitle{Напоминание из курса информатики :)}
\begin{itemize}
\item Всё, что умеет делать процессор -- операции над числами (арифметические, логические) и работать с периферийными устройствами (память, жёсткий диск, видеокарта, экран)
\pause
\begin{itemize}
\item Арифметическая операция: прочитать аргументы из памяти в процессор, выполнить операцию, записать результат в память
\pause
\item Вывести картинку: отправить цвета пикселей на экран
\pause
\item Записать файл: послать данные и их положение на диске в контроллер жесткого диска
\end{itemize}
\end{itemize}
\end{frame}

\begin{frame}[fragile]
\frametitle{Процессор}
\begin{itemize}
\item У процессора есть
\pause
\begin{itemize}
\item Набор поддерживаемых команд
\pause
\item Регистры -- ячейки памяти внутри процессора (очень быстрые чтение/запись)
\end{itemize}
\pause
\item Всё это вместо называется \textit{архитектурой процессора}: x86, x86-64 (AMD64), Aarch64 (ARM64), ...
\end{itemize}
\end{frame}

\begin{frame}[fragile]
\frametitle{Программа в машинных кодах}
\begin{itemize}
\item Как могла бы выглядеть (условно) программа (для гипотетической архитектуры процессора):
\pause
\usemintedstyle{lightbulb}
\begin{tcblisting}{listing engine=minted, minted language=cpp, minted options={fontsize=\scriptsize}, listing only, colback=codebg, colframe=codebg, rounded corners}
45 100 1    // Загрузить число из ячейки 100 в регистр R1
45 300 2    // Загрузить число из ячейки 300 в регистр R2
27 1 2 3    // Сложить R1 и R2 и записать результат в R3
46 500 3    // Записать число из R3 в ячейку памяти 500
\end{tcblisting}
\pause
\usemintedstyle{tango}
\item \mintinline{cpp}|27| -- код команды `сложить два числа'
\item \mintinline{cpp}|45| -- код команды `загрузить из памяти'
\item \mintinline{cpp}|46| -- код команды `записать в память'
\end{itemize}
\end{frame}

\begin{frame}[fragile]
\frametitle{Программа в машинных кодах}
\begin{itemize}
\item Программировать в таком виде очень неудобно:
\pause
\begin{itemize}
\item Нужно помнить номера команд и порядок их аргументов
\pause
\item Нужно помнить, какие регистры заняты (и чем)
\pause
\item Нужно помнить, какие ячейки памяти заняты (и чем)
\pause
\item Программа под одну архитектуру не будет работать на другой
\end{itemize}
\end{itemize}
\end{frame}

\begin{frame}[fragile]
\frametitle{Язык ассемблера}
\begin{itemize}
\item Ассемблер (от англ. assembler -- сборщик) -- программа, транслирующая текст на особом языке (языке ассемблера) в машинный код
\pause
\item У команд и регистров есть понятные, запоминающиеся имена
\pause
\item Ячейкам памяти можно давать свои имена
\pause
\usemintedstyle{lightbulb}
\begin{tcblisting}{listing engine=minted, minted language=cpp, minted options={fontsize=\scriptsize}, listing only, colback=codebg, colframe=codebg, rounded corners}
label 100 x
label 300 y
label 500 z

load x reg1
load y reg2
add reg1 reg2 reg3
store z reg3
\end{tcblisting}
\end{itemize}
\end{frame}

\begin{frame}[fragile]
\frametitle{Язык ассемблера}
\begin{itemize}
\item На языке ассемблера программировать гораздо удобнее, чем в машинных кодах, но некоторые проблемы остаются:
\pause
\begin{itemize}
\item Нет удобных управляющих конструкций (только jump'ы -- `прыжок' с одной инструкции на другую)
\pause
\item Тяжело писать сложные программы
\pause
\item Всё ещё нужно помнить, какие регистры чем заняты
\pause
\item Программа всё ещё заточена под конкретную архитектуру
\end{itemize}
\end{itemize}
\end{frame}

\begin{frame}[fragile]
\frametitle{Языки высокого уровня}
\usemintedstyle{lightbulb}
\begin{tcblisting}{listing engine=minted, minted language=cpp, minted options={fontsize=\scriptsize}, listing only, colback=codebg, colframe=codebg, rounded corners}
int x = 10;
int y = 20;
int z = x + y;
\end{tcblisting}
\pause
\begin{itemize}
\item Компилятор -- программа, преобразующая текст на языке программирования в язык ассемблера (и затем в машинный код)
\pause
\begin{itemize}
\item Сам подбирает нужные команды, ячейки памяти, номера регистров, и т.д.
\pause
\item Сам разбирается с особенностями конкретной архитектуры процессора
\pause
\item Оптимизирует подбор команд, чтобы программа работала быстрее
\end{itemize}
\end{itemize}
\end{frame}

\begin{frame}[fragile]
\frametitle{Язык C}
\begin{itemize}
\item Появился в 1972 в AT\&T Bell Labs с целью переписать ядро операционной системы Unix с языка ассемблера на кросс-платформенный язык
\pause
\item Был вдохновлён языками ALGOL 60, BCPL, B
\pause
\item Обрёл популярность после переписывания ядра Unix на C, и стал де-факто стандартом для низкоуровневых библиотек и интерфейсов
\pause
\item В современности, как правило, для поддержки новой архитектуры достаточно реализовать под неё компилятор C \textit{(или бэкенд LLVM)}
\end{itemize}
\end{frame}

\begin{frame}[fragile]
\frametitle{Программы из нескольких файлов}
\begin{itemize}
\item Зачем программу разбивают на части?
\pause
\begin{itemize}
\item Чем меньше код, тем его проще понять, тестировать, и отлаживать
\pause
\item Над разными частями могут работать разные люди
\pause
\item Выделенные части кода легче переиспользовать в других проектах
\end{itemize}
\end{itemize}
\end{frame}

\begin{frame}[fragile]
\frametitle{Пример программы: игра в шахматы}
\begin{itemize}
\item Файл с обработкой логики:
\begin{itemize}
\item Функция `проверить валидность хода'
\item Функция `сделать ход'
\item Функция `проверить выигрыш'
\end{itemize}
\pause
\item Файл с обработкой пользовательского ввода:
\begin{itemize}
\item Функция `узнать положение курсора мыши'
\item Функция `проверить нажатость кнопки мыши'
\end{itemize}
\pause
\item Файл с выводом на экран:
\begin{itemize}
\item Функция `нарисовать шахматную доску'
\item Функция `нарисовать шахматную фигуру'
\end{itemize}
\end{itemize}
\end{frame}

\begin{frame}[fragile]
\frametitle{Пример программы: сумма чисел}
\usemintedstyle{lightbulb}
\begin{tcblisting}{listing engine=minted, minted language=cpp, minted options={fontsize=\scriptsize}, listing only, colback=codebg, colframe=codebg, rounded corners}
// sum.c
int sum(int x, int y) {
  return x + y;
}
\end{tcblisting}
\usemintedstyle{lightbulb}
\begin{tcblisting}{listing engine=minted, minted language=cpp, minted options={fontsize=\scriptsize}, listing only, colback=codebg, colframe=codebg, rounded corners}
// main.c
int main() {
  int a = 3, b = 5;
  int c = sum(a, b);
  return 0;
}
\end{tcblisting}
\end{frame}

\begin{frame}[fragile]
\frametitle{Процесс сборки программы}
\begin{itemize}
\item Компилятор превращает исходный код в \textit{объектный файл} -- машинный код + пометки о том, где лежит код какой функции
\begin{itemize}
\item \mintinline{text}|main.c -> main.o|
\item \mintinline{text}| sum.c ->  sum.o|
\end{itemize}
\pause
\item \textit{Линковщик} (компоновщик) превращает список объектных файлов в готовую программу
\begin{itemize}
\item \mintinline{text}|main.o + sum.o -> sum.exe|
\end{itemize}
\end{itemize}
\end{frame}

\begin{frame}[fragile]
\frametitle{Линковщик}
\begin{itemize}
\item В объектных файлах вызовы функций записаны по именам, так как на этом этапе неизвестны финальные адреса функций в программе
\usemintedstyle{lightbulb}
\begin{tcblisting}{listing engine=minted, minted language=cpp, minted options={fontsize=\scriptsize}, listing only, colback=codebg, colframe=codebg, rounded corners}
[10] load a reg1
[11] load b reg2
[12] call sum
[13] store c reg3
[14] return
\end{tcblisting}
\end{itemize}
\end{frame}

\begin{frame}[fragile]
\frametitle{Линковщик}
\begin{itemize}
\item Линковщик собирает код всех функций вместе, назначает им адреса, и заменяет вызовы по имени на вызов по адресу
\usemintedstyle{lightbulb}
\begin{tcblisting}{listing engine=minted, minted language=cpp, minted options={fontsize=\scriptsize}, listing only, colback=codebg, colframe=codebg, rounded corners}
main:
[10] load a reg1
[11] load b reg2
[12] call [15]
[13] store c reg3
[14] return

sum:
[15] add reg1 reg2 reg3
[16] return
\end{tcblisting}
\pause
\item На уровне процессора имён функций \textit{не существует}, есть только адреса инструкций
\end{itemize}
\end{frame}

\begin{frame}[fragile]
\frametitle{Процесс сборки программы}
\begin{itemize}
\item Компиляция намного сложнее, чем линковка:
\pause
\begin{itemize}
\item Компиляция: превратить управляющие конструкции в jump'ы, проверить типы, выделить регистры, при необходимости выгрузить в память промежуточные вычисления, подобрать эффективные инструкции, удалить лишние инструкции
\pause
\item Линковка: склеить вместе готовый машинный код и заменить имена на адреса
\end{itemize}
\pause
\item Как правило, программист в конкретный момент работает только с небольшим набором файлов (например, изменили функцию \mintinline{cpp}|sum|, не меняя \mintinline{cpp}|main|)
\pause
\item Перекомпилировать нужно только те файлы, в которых были изменения (например, только \mintinline{cpp}|sum.c|)
\pause
\item Линковку всё равно придётся сделать целиком, но это \underline{дешёвая операция}
\end{itemize}
\end{frame}

\begin{frame}[fragile]
\frametitle{Ошибки линковки}
\begin{itemize}
\item \texttt{Undefined reference}: вызов несуществующей функции (вызов есть, а функции нет)
\item \texttt{Multiple definition}: несколько функций с одинаковым именем
\end{itemize}
\end{frame}

\begin{frame}[fragile]
\frametitle{Линковщик и типы}
\begin{itemize}
\item Линковщик и объектные файлы ничего не знают про типы и количество аргументов у функций
\pause
\usemintedstyle{lightbulb}
\begin{tcblisting}{listing engine=minted, minted language=cpp, minted options={fontsize=\scriptsize}, listing only, colback=codebg, colframe=codebg, rounded corners}
// sum.c
int sum(int x, int y) {
  return x + y;
}
\end{tcblisting}
\usemintedstyle{lightbulb}
\begin{tcblisting}{listing engine=minted, minted language=cpp, minted options={fontsize=\scriptsize}, listing only, colback=codebg, colframe=codebg, rounded corners}
// main.c
int main() {
  int a = 3, b = 5;
  int c = sum(a);
  return 0;
}
\end{tcblisting}
\pause
\usemintedstyle{tango}
\item Ошибка времени выполнения (runtime): фактически \mintinline{cpp}|sum| возьмёт что попало в качестве второго аргумента (то, что окажется в соответствующем регистре)
\end{itemize}
\end{frame}

\begin{frame}[fragile]
\frametitle{Заголовочные файлы}
\begin{itemize}
\item Для решения этой проблемы существуют \textit{заголовочные файлы (headers, хедеры)}
\pause
\item В специальном файле с расширением \mintinline{cpp}|.h| помещается \textit{объявление (declaration)} функции без её \textit{реализации (определение, definition)}
\usemintedstyle{lightbulb}
\begin{tcblisting}{listing engine=minted, minted language=cpp, minted options={fontsize=\scriptsize}, listing only, colback=codebg, colframe=codebg, rounded corners}
// sum.h
int sum(int x, int y);
\end{tcblisting}
\end{itemize}
\end{frame}

\begin{frame}[fragile]
\frametitle{Заголовочные файлы}
\begin{itemize}
\usemintedstyle{lightbulb}
\begin{tcblisting}{listing engine=minted, minted language=cpp, minted options={fontsize=\scriptsize}, listing only, colback=codebg, colframe=codebg, rounded corners}
// sum.c
#include "sum.h"

int sum(int x, int y) {
  return x + y;
}
\end{tcblisting}
\usemintedstyle{lightbulb}
\begin{tcblisting}{listing engine=minted, minted language=cpp, minted options={fontsize=\scriptsize}, listing only, colback=codebg, colframe=codebg, rounded corners}
// main.c
#include "sum.h"

int main() {
  int a = 3, b = 5;
  int c = sum(a);
  return 0;
}
\end{tcblisting}
\end{itemize}
\end{frame}

\begin{frame}[fragile]
\frametitle{Заголовочные файлы}
\begin{itemize}
\item \mintinline{cpp}|#include| -- \textit{директива препроцессора} (одной из частей компилятора C)
\pause
\item Она вставляет содержимое файла целиком в текущее место программы
\usemintedstyle{lightbulb}
\begin{tcblisting}{listing engine=minted, minted language=cpp, minted options={fontsize=\scriptsize}, listing only, colback=codebg, colframe=codebg, rounded corners}
// main.c после обработки препроцессором

int sum(int x, int y);

int main() {
  int a = 3, b = 5;
  int c = sum(a); // ошибка: неверное число аргументов
  return 0;
}
\end{tcblisting}
\pause
\item Имея \textit{объявление функции}, компилятор может проверить соответствие типов аргументов и их количества
\end{itemize}
\end{frame}

\end{document}
